% ============================================================
%  Conclusion
% ============================================================

\chapter*{Conclusion}
\addcontentsline{toc}{chapter}{Conclusion}
\markboth{Conclusion}{}

Ce projet a permis de reproduire, enrichir et améliorer le modèle MCTNet pour
la classification de cultures agricoles à partir de séries temporelles Sentinel-2,
à travers trois axes complémentaires.

\textbf{Partie 1.}
La reproduction de MCTNet a été menée de bout en bout : collecte GEE sur 8 zones
Arkansas et 4 zones Californie, prétraitement (masque données manquantes, normalisation,
reshape), et entraînement avec le protocole original. Les performances reproduites
(OA\,=\,0.9603 / Kappa\,=\,0.9504 en Arkansas ; OA\,=\,0.9311 / Kappa\,=\,0.9166
en Californie) confirment la reproductibilité du modèle et servent de baseline
pour toutes les contributions suivantes.

\textbf{Partie 2.}
L'intégration de covariables environnementales (sol, climat ERA5, topographie)
dans une variante MCTNetWithCovars apporte des gains substantiels en Arkansas
(jusqu'à +1.72 pts OA avec les seules propriétés pédologiques), région où les
cultures sont fortement contraintes par les conditions des plaines alluviales.
En Californie, seule la combinaison complète des trois groupes dépasse clairement
le baseline (+0.46 pts OA), reflétant la plus grande diversité géographique
et la présence de cultures pérennes.

\textbf{Partie 3.}
Trois contributions architecturales ont été évaluées :
\begin{itemize}
  \item \textbf{GatedMCTNet} améliore la baseline de +0.92 pts OA en Arkansas
        avec seulement 4\,270 paramètres supplémentaires (+7.5\%). En Californie,
        la gate nécessite davantage de données pour converger, entraînant une
        légère régression.
  \item \textbf{MCTNet Multiscale} apporte des gains modestes mais cohérents
        sur les deux régions (+0.65 pts F1 Arkansas, +0.10 pts F1 Californie),
        au prix d'une multiplication par 2.75 du nombre de paramètres
        (56\,798 $\to$ 156\,198), avec un risque accru de surapprentissage.
  \item \textbf{MCTNetUSkip} réintègre les représentations intermédiaires des
        stages 1 et 2 via des connexions skip légères, améliorant les résultats
        en Arkansas (+0.34 pts OA) avec \emph{moins} de paramètres que MCTNet.
        Son extension en encodeur-décodeur U-Net avec covariables atteint
        OA\,=\,0.9785 en Arkansas (configuration Sol+Topo) et OA\,=\,0.9395
        en Californie (configuration complète).
\end{itemize}

Ces travaux ouvrent plusieurs perspectives : consolidation des résultats de
GatedMCTNet en Californie sur un jeu d'entraînement plus large, régularisation
du Multiscale (dropout, pruning des kernels redondants), et exploration de
mécanismes d'attention croisée entre la branche sérielle et la branche
covariables pour mieux exploiter l'information exogène.
